\paragraph{Problem 3 answer}\GradescopeComment{This should be on page 10 (top) of your PDF.}

\begin{enumerate}[label=(\alph*)]

  \item   % Per Section 7.4, define the subproblems your algorithm solves. Which is the final answer?

    Here is a definition of all subproblems that the algorithm solves, given $(a_1, a_2, \ldots, a_n)$.

    \begin{blue}
      \REPLACEME
    \end{blue}

    The final answer is \BLUE{\( \REPLACEME \)}.

  \item % State an appropriate recurrence relation, including boundary cases.

    Here is the recurrence relation:

    \begin{blue}
      \[\REPLACEME =
        \begin{cases}
          \REPLACEME
        \end{cases}
      \]
    \end{blue}
    
    % For reference, here is the LaTeX source for the recurrence for longest palindromic subsequence, Section 7.10.
    % \[M(i, j)  =
    %   \begin{cases}
    %     0 & \text{if } j < i, \\
    %     1 & \text{if } j = i, \\
    %     \max\big(M(i, j-1), M(i+1, j), 2 + M(i+1, j-1)\big) & \text{if } j>i \text{ and } X[i] = X[j], \\
    %     \max\big(M(i, j-1), M(i+1, j)\big) & \text{if } j>i \text{ and } X[i] \ne X[j].
    %   \end{cases}
    % \]

\item % Explain clearly in at most five lines why the recurrence relation is correct.

  Here is why the recurrence relation is correct:

  \begin{blue}
    \REPLACEME
  \end{blue}
  
\end{enumerate}

\newpage

\begin{enumerate}

\item[(d)] % Give a tight bound on the big-$O$ running time of your algorithm in terms of $n$.

  The big-$\Theta$ running time of the resulting algorithm is \GradescopeComment{This should be on page 11 (top) of your PDF.}
  \BLUE{
    \[
      \Theta(\REPLACEME).
    \]
  }

\item[(e)] % Explain clearly in just 1--3 lines why the bound holds.

  Here is why the bound above holds:

  \begin{blue}
    \REPLACEME
  \end{blue}
  
\end{enumerate}

\begin{itemize}
  
\item[(f)] % Give pseudo-code for an iterative implementation of the resulting algorithm.
  Here is pseudo-code for the iterative implementation of the resulting algorithm.
  
  \begin{blue}
    \REPLACEME
  \end{blue}
  
\item[(g)]
  % Submit your algorithm to the Hacker-Rank challenge at the URL in the footnote below.\footnote
  % {\url{https://www.hackerrank.com/nealeyoung-algs101-game-of-piles}}
  % What is the username of your Hacker-Rank submission?
  % For full credit, your submission should pass all tests.

  The username for my Hacker-Rank submission is \BLUE{
    \REPLACEME
  }.

  % NOTE: TO REPRESENT AN UNDERSCORE USE  \_
  
\end{itemize}


\vfill

\paragraph{Problem \arabic{problem} collaboration and citations}~\GradescopeComment{This should be on page 11 (bottom) of your PDF.}

I collaborated on this problem with the following people:
\begin{blue}
  \begin{itemize}
  \item \REPLACEME                    % write "none" if none
  \end{itemize}
\end{blue}

My answers use ideas from the following sources (other than the lecture notes and textbooks): 
\begin{blue}
  \begin{itemize}
  \item \REPLACEME                    % write "none" if none
  \end{itemize}
\end{blue}



%%% Local Variables:
%%% mode: latex
%%% TeX-master: "homework_7"
%%% End:
